% RESEARCH_DRAFT_SHA256: 4b2ccf18ea8da9746025eb01fd07f09e5cff1b3e27b74cca8f6ce76c07f4f24f
% Layout and source-reference wrappers expanded for the web renderer.
\documentclass[11pt]{article}
\usepackage[margin=25mm]{geometry}
\usepackage{fontspec,unicode-math}
\setmainfont{Latin Modern Roman}
\setmathfont{Latin Modern Math}
\usepackage{amsmath,amssymb,mathtools,xurl,hyperref}
\hypersetup{colorlinks=true,urlcolor=blue,linkcolor=blue}
\setlength{\parindent}{0pt}
\setlength{\parskip}{5pt}
\setlength{\emergencystretch}{3em}
\newcommand{\sref}[2]{\hyperlink{source:#1:#2}{[S#2]}}
\newcommand{\cataloguescope}[1]{\paragraph{Source question.}#1}
\newcommand{\floor}[1]{\left\lfloor#1\right\rfloor}
\begin{document}

\section*{Duplicate subtree-hull record: a gluing attempt and its distributional gap}
\noindent UnsolvedMath ID 30001103; OWR-2489-009. Research attempt dated 11 October 2026.
\paragraph{Result in scope.}This manuscript contains a source-specific attempt and complete proofs of the stated auxiliary calculations. It makes no new solution or novelty claim for the entire target.
\section{Original problem and preserved definition}
The following definition and source list are reproduced from the repository. Any earlier inspection dates and reported literature claims in this preserved text describe that source record, not additional verification by this attempt.
% BEGIN PRESERVED DEFINITION
\subsection{Definitions and mathematical statement}

\paragraph{Shared notation.}Standard number systems are $\mathbb N,\mathbb Z,\mathbb Q,\mathbb R,\mathbb C$. Specialized notation and index ranges follow the cited source and the definitions below.


For a subtree $T$, its constraints and integral coordinates are restricted from the graph. The conjectured equality compares the convex hull of all mixed-integer solutions with the simultaneous convex-hull restrictions from the specified subtrees.
\paragraph{Record scope and qualifications.}
The catalogue title is retained as metadata. Named certificates, contracts, bounds or frameworks in that title are not asserted as results of the original source extract.
The supplied registry flags this extraction as unrecoverable. It is displayed as a literal source fragment; the definitions provide context rather than a replacement formal statement.

Source-record limitation: Source flags the statement as unrecoverable: The official report confirms that this repeats the same Conjecture 2 extracted as OWR-2489-004. No separate clean statement is promoted; the original record remains publicly visible for provenance.

Related catalogue records: \href{https://mehmetmars7.github.io/Erdosproblems-llm-hunter/problem.html?type=open_problems&id=30001100}{30001100}. Separate source IDs are preserved for their provenance.
The source text below is rendered literally where necessary; damaged formulas and composite extracts are not silently repaired into a new mathematical claim.
The supplied catalogue marks this record as unrecoverable. Its extract and terminology are documented, but a complete original problem has not been reconstructed.
\paragraph{Statement recorded in the supplied source.}
With S(G,I) as defined in the report and T the subtrees of G whose integral vertices are exactly the leaves, is the convex-hull identity conv(S(G,I))=\textbackslash{}bigcap\_\{T\}conv(S(T,I\_T)) valid for every denominator parameter k? It is known for k=2 and remains open for every k\textbackslash{}geq3. \href{https://doi.org/10.4171/owr/2008/51}{[S2]}
\paragraph{Source status.}
The supplied catalogue labels this source record open; that metadata was checked against the supplied snapshot on 2026-10-11. The label is preserved for provenance and is not an independent verification of present open status or mathematical correctness. The catalogue literature assessment is dated 2026-08-21.
\subsection{Short English statement}
The supplied fragment concerns subtree Decomposition of Mixed-Integer Convex Hulls. Its complete mathematical target is not recovered.
\subsection{Sources}
\begin{itemize}
\item[S1] Ulam AI and the UnsolvedMath Contributors, UnsolvedMath: A Curated Collection of Open Mathematics Problems, record 30001103 (OWR-2489-009). Catalogue curation and metadata: CC BY 4.0. The source extract is retained; definitions, English summaries and qualifications are editorial additions.
\url{https://huggingface.co/datasets/ulamai/UnsolvedMath}.
\item[S2] Combinatorial Optimization (2009), Oberwolfach Reports, DOI 10.4171/owr/2008/51, pp. 2875--2942
\url{https://doi.org/10.4171/owr/2008/51}.
\item[S3] network formulations
\url{https://doi.org/10.1287/moor.1080.0354}.
\item Source-specific Oberwolfach Reports licences, including CC BY-SA 4.0 where applicable, are retained; see THIRD\_PARTY\_NOTICES.md.
\end{itemize}
% END PRESERVED DEFINITION
\subsection{Verbatim supplied catalogue question}
With S(G,I) as defined in the report and T the subtrees of G whose integral vertices are exactly the leaves, is the convex-hull identity conv(S(G,I))=\textbackslash{}bigcap\_\{T\}conv(S(T,I\_T)) valid for every denominator parameter k? It is known for k=2 and remains open for every k\textbackslash{}geq3.

\section{Preserved duplicate and source scope}
The supplied record is flagged unrecoverable because it repeats the same
Conforti Conjecture 2 as record 30001100, rather than defining a second
independent problem. Inspection of the primary report confirms the shared
subtree identity. The denominator $k$ is the least positive integer making
all rational edge right-hand sides integral. This manuscript retains the
flag and original fragment and investigates a gluing proof of that shared
question. It makes no claim to reconstruct a different missing theorem.

\paragraph{The literal ambient-domain omission persists.}
On a single edge $uv$ with $I=\{u\}$ and $b_{uv}=1/3$, the edge-tree
has two leaves but only one integral vertex, so it is excluded from the
specified family. Any included singleton subtree has hull equal to its
whole coordinate line, since both $\operatorname{conv}(\mathbb Z)$ and
$\mathbb R$ equal $\mathbb R$. Thus the lifted subtree intersection imposes
no constraint, while the global hull satisfies $x_u+x_v\ge1/3$.
This also happens at denominator two, contradicting the source's stated
known case if all hulls are interpreted in unrestricted ambient space.
An implicit continuous relaxation or different subtree convention is therefore
missing. This diagnostic is not treated as a resolution of the intended
conjecture. All subsequent gluing observations concern the proposed proof
mechanism, not a silent repair of the original statement.

\section{Attempt to glue convex decompositions}
Membership of $\bar x_T$ in a subtree hull supplies a finite distribution
on feasible local points whose coordinate mean is $\bar x_T$. A tempting
proof would glue these local distributions into a globally feasible random
point with mean $\bar x$. Means on common coordinates are not enough to
justify that step.

\paragraph{A precise gluing lemma under a stronger hypothesis.}
Let finite distributions $P_A,P_B$ on coordinates indexed by $A,B$ agree
as entire marginal distributions on $S=A\cap B$, with common marginal
$\mu$. For consistent assignments $a,b$ having common restriction $s$
with $\mu(s)>0$, define
\[
P(a,b)=\frac{P_A(a)P_B(b)}{\mu(s)};
\]
put zero mass on all other pairs. Summing over compatible $b$ gives
$P_A(a)\mu(s)/\mu(s)=P_A(a)$; if $\mu(s)=0$, every assignment
restricting to $s$ has zero marginal mass, so nothing is lost.
The same argument gives marginal $P_B$, and summing either marginal proves
normalization. Every positive-mass global assignment has allowed local
restrictions. Hence a gluing is possible with those full marginal hypotheses.
The same construction can be iterated along a tree of variable blocks when,
at each attachment, the intersection with all earlier blocks is the specified
separator and its full marginal agrees. This follows by induction using the
two-block argument, which preserves the already constructed marginal exactly.

\paragraph{Equal first moments do not imply compatible local laws.}
A single coordinate with law supported equally on $\{-1,1\}$ has the
same mean zero as the point mass at zero, but these distributions cannot be
the two marginals of one shared coordinate. Thus even two local
convex-combination witnesses cannot simply be combined from their means.
Moreover cycle consistency can fail despite agreement of all singleton
marginals. On four binary variables consider the relations
\[
x_1=x_2,\quad x_2=x_3,\quad x_3=x_4,\quad x_4=1-x_1.
\]
Each individual relation admits a distribution on its two allowed pairs
with each coordinate fair. More strongly, every forest obtained by deleting
at least one cycle edge admits two complementary assignments on each
component: choose a root bit and propagate equality or complementation
along the unique paths. Fairly averaging complementary assignments gives
mean $(1/2,1/2,1/2,1/2)$ on every such forest.
Yet no global assignment exists, since the first three equalities imply
$x_4=x_1$, contradicting the last relation. The incidence graph is an
even cycle, so bipartiteness alone cannot make this generic gluing inference
valid. These relations are not the source's one-sided covering inequalities;
this is explicitly a counterexample to a general proof principle, not to
the intended mixed-integer covering conjecture.

\section{Adversarial audit and gap}
Global feasible points always restrict to locally feasible points, and their
finite convex combinations do too, proving the easy global-to-local inclusion.
The reverse needs a structural argument producing compatible distributions,
or another way to use the covering inequalities; the gluing lemma does not
supply that compatibility from means. The audit separated full distributions
from first moments, acyclic attachment from cycle consistency, and the generic
binary counterexample from the exact source class. The duplicate status and
literal-domain defect remain visible. No computation was performed and no
independent target or new result is counted from this duplicate record.

\section{Process and contributions}
Mathematics and model-assisted adversarial self-audit: OpenAI gpt-6-astra, ultra reasoning, as dispatched for this research batch; exact serving revision not exposed. Codex (GPT-6; exact serving revision not exposed) materially assisted with orchestration and format review. No human mathematical authorship, mathematical edits, or independent human proof checking is asserted.
Contribution prompt A was applied to the full definition and supplied source row, then prompt B was applied adversarially to the derived results. The item-specific audit above records the scope checks and rejected extrapolations. Local assembly and structural checks do not establish mathematical correctness. The source catalogue status is preserved in TOP\_PROBLEM and is separate from this attempt status.
\paragraph{Integration review.} Codex (GPT-6; exact serving revision not exposed) reviewed the written arguments and expanded layout and source-reference wrappers. Repository metadata and the production web renderer were checked using Python 3.13 and JavaScript with MathJax 3.2.2. This model-assisted review is not independent human proof verification; no PDF compilation or formal proof checking is asserted.

\section{Completion Estimate}
\noindent\textbf{COMPLETION: 0\%}
The proved material is a rederivation of known or auxiliary facts rather than a substantiated new advance on the original target. The shared intended conjecture remains unresolved: local hull memberships provide means, not the compatible full marginal distributions needed by the gluing construction; the source fragment also omits an ambient convention.
\section{References}
The complete original Sources subsection above is retained, including catalogue attribution and licence notices. The additional list records primary sources actually consulted for this attempt on 11 October 2026; uncited source links reproduced above were not all rechecked.
\begin{itemize}
\item Oberwolfach Report 51/2008: Conforti, Combinatorial Mixed-Integer Programming, printed pp.2904--2905; Alcuin-number contribution, pp.2891--2893. Each relevant original passage read from cached primary PDF text on 11 October 2026. \url{https://ems.press/content/serial-article-files/46195}.
\end{itemize}
\section{Final status}
\textbf{Unresolved in this attempt.} The shared intended conjecture remains unresolved: local hull memberships provide means, not the compatible full marginal distributions needed by the gluing construction; the source fragment also omits an ambient convention.
\end{document}
